\documentclass[11pt]{article}
\usepackage[margin=2.1cm]{geometry}
\usepackage{booktabs}
\usepackage{amsmath}
\usepackage{enumitem}
\usepackage[T1]{fontenc}
\usepackage{lmodern}

\setlength{\parindent}{0pt}
\setlength{\parskip}{5pt}
\newcommand{\code}[1]{\texttt{#1}}
\newcommand{\answerbox}[1]{\vspace{2pt}\fbox{\parbox[c][#1][c]{\dimexpr\linewidth-2\fboxsep-2\fboxrule}{\ }}\vspace{4pt}}

\begin{document}

\begin{center}
{\Large\bfseries EE\,362 --- Pre-Lab 8}\\[2pt]
{\large Next Instruction Address Generation II}\\[6pt]
\begin{tabular}{@{}p{6.5cm}p{4cm}p{4cm}@{}}
Name: \hrulefill & Section: Mon / Tue & Date: \hrulefill
\end{tabular}
\end{center}

\vspace{2pt}\hrule\vspace{4pt}

{\small\textbf{Due at the start of the lab session.} Five points, about twenty
minutes, on paper.}

\vspace{4pt}\hrule

\section*{Question 1 --- The six equations}

\textbf{(a)} The ALU gives you two flags: $Z$ (result is zero) and $N$ (result
is negative). Write the Boolean equation for \emph{Taken} for each condition.

\begin{center}
\begin{tabular}{lll@{\hskip 8pt}||@{\hskip 8pt}l}
\toprule
code & meaning & & Taken $=$ \\
\midrule
\code{NONE} & always   & & \rule{5cm}{0.4pt} \\[5pt]
\code{EQ}   & $=0$     & & \rule{5cm}{0.4pt} \\[5pt]
\code{NEQ}  & $\neq 0$ & & \rule{5cm}{0.4pt} \\[5pt]
\code{GTZ}  & $>0$     & & \rule{5cm}{0.4pt} \\[5pt]
\code{GEZ}  & $\geq 0$ & & \rule{5cm}{0.4pt} \\[5pt]
\code{LTZ}  & $<0$     & & \rule{5cm}{0.4pt} \\[5pt]
\code{LEZ}  & $\leq 0$ & & \rule{5cm}{0.4pt} \\
\bottomrule
\end{tabular}
\end{center}

\textbf{(b)} Three of these pairs are complements of each other. Circle them in
the table above, and say in one sentence what that tells you about how much
hardware the seven conditions actually cost.

\answerbox{1.6cm}

\textbf{(c)} A classmate writes \code{GTZ} as $\overline{N}$. Give the exact
flag values where their version and the correct one disagree, and say which
instruction would then behave wrongly.

\answerbox{2.0cm}

\section*{Question 2 --- Absolute jumps}

Here is a real instruction from the Lab 12 program:

\begin{center}
\begin{tabular}{lll}
\toprule
\code{0x00000014} & \code{0x08000002} & \code{J 0x08} \\
\bottomrule
\end{tabular}
\end{center}

\textbf{(a)} Compute the target. Show the three pieces.

\vspace{4pt}
\code{NSA} $=$ \rule{3cm}{0.4pt} \hfill
top 4 bits $=$ \rule{2.5cm}{0.4pt} \hfill
\code{Abslt} $=$ \rule{4cm}{0.4pt}
\vspace{8pt}

\textbf{(b)} $4 + 26 + 2 = 32$. Say where each group of bits comes from, and
then state the range limit this produces --- how far can a \code{J} reach, and
why?

\answerbox{2.4cm}

\textbf{(c)} What instruction do you use when you need to go further than that,
and what makes it able to?

\answerbox{1.4cm}

\section*{Question 3 --- Test plan}

\textbf{(a)} \code{Sel} and \code{Cond} are two separate control fields. Give
one pair of test vectors that differ \emph{only} in \code{Taken} and produce
two different next addresses --- and one pair that differ only in \code{Taken}
and produce the \emph{same} address.

\answerbox{2.6cm}

\textbf{(b)} Sketch your scripted sequence: the list of steps the PC will walk
through, and the address you expect after each one.

\begin{center}
\begin{tabular}{@{}c@{\hskip 8pt}p{5.4cm}@{\hskip 8pt}p{3cm}@{\hskip 8pt}p{3.4cm}@{}}
\toprule
\# & what happens & expected PC & what it proves \\
\midrule
1 & \rule{5.2cm}{0.4pt} & \rule{2.8cm}{0.4pt} & \rule{3.2cm}{0.4pt} \\[7pt]
2 & \rule{5.2cm}{0.4pt} & \rule{2.8cm}{0.4pt} & \rule{3.2cm}{0.4pt} \\[7pt]
3 & \rule{5.2cm}{0.4pt} & \rule{2.8cm}{0.4pt} & \rule{3.2cm}{0.4pt} \\[7pt]
4 & \rule{5.2cm}{0.4pt} & \rule{2.8cm}{0.4pt} & \rule{3.2cm}{0.4pt} \\[7pt]
5 & \rule{5.2cm}{0.4pt} & \rule{2.8cm}{0.4pt} & \rule{3.2cm}{0.4pt} \\[7pt]
6 & \rule{5.2cm}{0.4pt} & \rule{2.8cm}{0.4pt} & \rule{3.2cm}{0.4pt} \\
\bottomrule
\end{tabular}
\end{center}

\vfill
{\small\emph{Question 3 is the one that matters.} It is the same question as
Section 5 of your lab report.}

\end{document}
